\documentclass[11pt,a4paper]{article}

% ============================================================================
% Packages
% ============================================================================
\usepackage[margin=2.5cm]{geometry}
\usepackage{amsmath,amssymb}
\usepackage{graphicx}
\usepackage{booktabs}
\usepackage{longtable}
\usepackage{multirow}
\usepackage{float}
\usepackage[dvipsnames]{xcolor}
\usepackage{caption}
\usepackage{subcaption}
\usepackage{hyperref}
\usepackage{enumitem}
\usepackage{siunitx}
\usepackage{bm}

\hypersetup{
  colorlinks=true,
  linkcolor=NavyBlue,
  citecolor=OliveGreen,
  urlcolor=BrickRed,
}

\sisetup{
  per-mode = symbol,
  inter-unit-product = \ensuremath{{}\cdot{}},
}

% ============================================================================
% Custom commands
% ============================================================================
\newcommand{\species}[1]{\ensuremath{\mathrm{#1}}}
\newcommand{\ee}{\species{e}}
\newcommand{\ddt}[1]{\frac{\partial #1}{\partial t}}
\newcommand{\dv}[2]{\frac{d #1}{d #2}}
\newcommand{\Te}{\ensuremath{T_e}}
\newcommand{\Tg}{\ensuremath{T_g}}
\renewcommand{\ne}{\ensuremath{n_e}}
\newcommand{\mev}{\,\mathrm{eV}}

\title{\bfseries A Zero-Dimensional Global Model for Ar/NF$_3$\\
  Inductively Coupled Plasmas:\\
  Mathematical Formulation, Implementation, and Results}
\author{Computational Plasma Physics --- Model Reconstruction Report}
\date{\today}

\begin{document}
\maketitle

\begin{abstract}
  A zero-dimensional (0D) global model is developed for an Ar/NF$_3$ inductively
  coupled plasma (ICP) discharge, reproducing the framework of Huang \textit{et al.}\
  (2026). The model solves time-dependent particle balance equations for 21 plasma
  species coupled with an electron energy balance, incorporating 155 volumetric
  reactions (electron-impact, Penning, neutral--neutral, ion--neutral, ion--ion,
  thermal dissociation, and three-body recombination) together with surface
  recombination losses and gas pumping.  The ion wall-loss formulation follows the
  Thorsteinsson--Gudmundsson approach for electronegative plasmas.  This report
  presents the complete mathematical framework, all governing equations, the
  numerical implementation in Python, and parametric results for power and pressure
  sweeps at 20\% NF$_3$ feed composition.  The model correctly reproduces the key
  qualitative trends: a linear increase of electron density with absorbed power, a
  monotonic decrease of electron density and electron temperature with pressure, near-complete
  dissociation of NF$_3$, Ar$^+$ as the dominant positive ion, and F$^-$ as the
  dominant negative ion.
\end{abstract}

\tableofcontents
\newpage

% ============================================================================
\section{Introduction}
\label{sec:intro}
% ============================================================================

Nitrogen trifluoride (NF$_3$) is widely used as a remote plasma cleaning gas in
semiconductor manufacturing, particularly for chamber cleaning in chemical vapor
deposition (CVD) and atomic layer deposition (ALD) tools. Understanding the
dissociation chemistry of NF$_3$ in argon carrier gas is critical for optimizing
fluorine atom delivery and minimizing greenhouse gas emissions.

This report documents the reconstruction of a 0D global (volume-averaged) model
for an Ar/NF$_3$ ICP reactor, following the formulation of Huang \textit{et al.}\
(\textit{Plasma Sources Sci.\ Technol.}\ \textbf{35}, 015019, 2026). The model
tracks 21 species through 155 gas-phase reactions and includes surface
recombination, gas pumping, and a self-consistent electron energy equation.

The primary objectives are:
\begin{enumerate}[nosep]
  \item Formulate the complete mathematical model.
  \item Implement the model in Python using an ODE integrator.
  \item Reproduce the parametric trends reported in the reference paper.
  \item Identify and document all assumptions and approximations.
\end{enumerate}


% ============================================================================
\section{Species and Reaction Mechanism}
\label{sec:species}
% ============================================================================

\subsection{Species Set}

The model tracks 21 distinct species in four categories:

\begin{table}[H]
  \centering
  \caption{Plasma species included in the global model.}
  \label{tab:species}
  \begin{tabular}{lll}
    \toprule
    \textbf{Category} & \textbf{Count} & \textbf{Species} \\
    \midrule
    Neutral ground states & 8 &
      NF$_3$, NF$_2$, NF, F$_2$, N$_2$, F, N, Ar \\
    Neutral excited states & 2 &
      Ar($1s_5$), Ar($4p$) \\
    Positive ions & 8 &
      NF$_3^+$, NF$_2^+$, NF$^+$, F$_2^+$, N$_2^+$, F$^+$, N$^+$, Ar$^+$ \\
    Negative ions & 2 &
      F$^-$, F$_2^-$ \\
    \midrule
    Electrons & 1 & e \\
    \bottomrule
  \end{tabular}
\end{table}

The electron density is not solved as an independent ODE variable; instead it is
obtained at every time step from the charge neutrality constraint (Section~\ref{sec:neutrality}).

\subsection{Reaction Mechanism}

A total of 155 reactions are included, organized into 10 categories
(Table~\ref{tab:reactions_summary}).

\begin{table}[H]
  \centering
  \caption{Summary of the 155 gas-phase and surface reactions.}
  \label{tab:reactions_summary}
  \begin{tabular}{llr}
    \toprule
    \textbf{Category} & \textbf{Indices} & \textbf{Count} \\
    \midrule
    Electron--NF$_x$ (ionization, dissociation, attachment, excitation)
      & R1--R36 & 36 \\
    Electron--F$_x$ (ionization, dissociation, attachment, excitation)
      & R37--R49 & 13 \\
    Electron--N$_x$ (ionization, dissociation, excitation)
      & R50--R65 & 16 \\
    Electron--Ar (ionization, excitation, superelastic)
      & R66--R77 & 12 \\
    Electron--ion recombination and detachment
      & R78--R86 & 9 \\
    Penning dissociation (Ar$^*$ + NF$_x$)
      & R87--R89 & 3 \\
    Penning ionization (Ar$^*$ + NF$_x$, Ar$^*$)
      & R90--R93 & 4 \\
    Neutral--neutral reactions
      & R94--R102 & 9 \\
    Ion--neutral charge exchange
      & R103--R136 & 34 \\
    Ion--ion neutralization
      & R137--R145 & 9 \\
    High-temperature thermal dissociation
      & R146--R150 & 5 \\
    Three-body recombination
      & R151--R155 & 5 \\
    \midrule
    \textbf{Total} & & \textbf{155} \\
    \bottomrule
  \end{tabular}
\end{table}


\subsection{Rate Coefficient Parameterization}
\label{sec:rate_coeff}

All electron-impact rate coefficients are parameterized in modified Arrhenius form:
\begin{equation}
  \label{eq:arrhenius_Te}
  k(\Te) \;=\; A\;\Te^{\,B}\;\exp\!\Bigl(-\frac{C}{\Te^{\,D}}\Bigr),
\end{equation}
where $\Te$ is the electron temperature in eV, and $A$, $B$, $C$, $D$ are
tabulated constants for each reaction. For a subset of reactions involving N$_2$
excitation (R54--R57), the rate coefficients are parameterized with $\Te$ in Kelvin
($\Te[\mathrm{K}] = 11604.5\;\Te[\mathrm{eV}]$) and must not be unit-converted before
evaluation.

Heavy-particle reaction rates (neutral--neutral, ion--neutral, three-body) use a
gas-temperature-dependent Arrhenius form:
\begin{equation}
  \label{eq:arrhenius_Tg}
  k(\Tg) \;=\; A\;\Bigl(\frac{\Tg}{300}\Bigr)^{\!B}\;
  \exp\!\Bigl(-\frac{E_a}{\Tg}\Bigr),
\end{equation}
where $\Tg$ is the gas temperature in Kelvin and $E_a$ is an activation energy in~K.


% ============================================================================
\section{Governing Equations}
\label{sec:equations}
% ============================================================================

\subsection{Particle Balance}

For each heavy species $X$ (neutrals and ions), the time-dependent particle
balance is:
\begin{equation}
  \label{eq:particle_balance}
  \boxed{
  \ddt{n_X} \;=\;
    \underbrace{\sum_r \nu_{X,r}\,R_r}_{\text{volumetric}}
    \;-\; \underbrace{k_{\text{wall},X}\,n_X}_{\text{wall loss}}
    \;+\; \underbrace{\dot{n}_{X,\text{pump}}}_{\text{pumping}}
  }
\end{equation}
where $n_X$ is the number density (m$^{-3}$), $\nu_{X,r}$ is the net
stoichiometric coefficient of species $X$ in reaction $r$ (positive for
production, negative for consumption), $R_r$ is the volumetric reaction rate
(m$^{-3}$\,s$^{-1}$), and the wall loss and pumping terms are detailed below.

The volumetric rate of a two-body electron-impact reaction $\ee + X \to
\text{products}$ is:
\begin{equation}
  R_r = k_r(\Te)\;\ne\;n_X,
\end{equation}
while for a heavy-particle two-body reaction $A + B \to \text{products}$:
\begin{equation}
  R_r = k_r(\Tg)\;n_A\;n_B.
\end{equation}
For three-body reactions $A + B + M \to \text{products}$:
\begin{equation}
  R_r = k_r(\Tg)\;n_A\;n_B\;n_{\text{total}},
\end{equation}
where $n_{\text{total}} = \sum_i n_i$ over all neutral species.


\subsection{Charge Neutrality}
\label{sec:neutrality}

Rather than solving an ODE for $\ne$, the electron density is determined
algebraically at each time step:
\begin{equation}
  \label{eq:neutrality}
  \boxed{
  \ne = \sum_{i \in \text{cations}} n_i \;-\; \sum_{j \in \text{anions}} n_j
  }
\end{equation}
In this model, the cations are \{NF$_3^+$, NF$_2^+$, NF$^+$, F$_2^+$,
N$_2^+$, F$^+$, N$^+$, Ar$^+$\} and the anions are \{F$^-$, F$_2^-$\}.


\subsection{Neutral Wall Loss}
\label{sec:neutral_wall}

The loss of a reactive neutral species $X$ at the chamber walls with surface
recombination probability $\gamma_X$ is modeled as a first-order loss:
\begin{equation}
  \label{eq:neutral_wall}
  k_{\text{wall},X} = \left[
    \frac{\Lambda^2}{D_X} + \frac{2V}{A\,\bar{v}_X}\,\frac{2-\gamma_X}{\gamma_X}
  \right]^{-1},
\end{equation}
where $V$ is the chamber volume, $A$ is the total wall area, and:

\paragraph{Diffusion length:}
\begin{equation}
  \label{eq:diff_length}
  \frac{1}{\Lambda^2} = \left(\frac{\pi}{L}\right)^{\!2}
    + \left(\frac{2.405}{R}\right)^{\!2},
\end{equation}
with $R$ the chamber radius and $L$ the chamber height.

\paragraph{Free diffusion coefficient:}
\begin{equation}
  \label{eq:diffusion}
  D_X = \frac{k_B\Tg\,\lambda_X}{m_X\,\bar{v}_X},
\end{equation}
where $\lambda_X = 1/(n_{\text{total}}\,\sigma_{nn})$ is the mean free path
with $\sigma_{nn} \approx \SI{3e-19}{m^2}$ the neutral--neutral collision cross
section.

\paragraph{Thermal velocity:}
\begin{equation}
  \label{eq:vth}
  \bar{v}_X = \sqrt{\frac{8\,k_B\Tg}{\pi\,m_X}}.
\end{equation}

The surface recombination probabilities used are: $\gamma_\text{F} = 0.01$,
$\gamma_\text{N} = 0.07$, $\gamma_\text{NF}_x = 0.01$, $\gamma_{\text{Ar}^*}
= 1.0$. Upon recombination at the wall, atoms form their parent molecules
(F~$\to$~$\frac{1}{2}$\,F$_2$, N~$\to$~$\frac{1}{2}$\,N$_2$) and excited
argon de-excites to the ground state.


\subsection{Ion Wall Loss}
\label{sec:ion_wall}

The positive ion wall-loss rate follows the Thorsteinsson--Gudmundsson
formulation for electronegative plasmas:
\begin{equation}
  \label{eq:ion_wall}
  \boxed{
  \left(\ddt{n_i}\right)_{\!\text{wall}} = -\frac{u_{B,i}\,A_{\text{eff}}}{V}\;n_i
  }
\end{equation}

\paragraph{Modified Bohm velocity:}
\begin{equation}
  \label{eq:bohm}
  u_{B,i} = \sqrt{\frac{e\Te}{m_i}\;\frac{1+\alpha_s}{1+\gamma_-\,\alpha_s}},
\end{equation}
where $\alpha_s = n_-^{(s)}/\ne^{(s)}$ is the electronegativity at the sheath
edge, $\gamma_- = \Te/T_-$, and $m_i$ is the ion mass.

\paragraph{Effective wall area:}
\begin{equation}
  \label{eq:Aeff}
  A_{\text{eff}} = 2\pi\bigl(R^2\,h_L + R\,L\,h_R\bigr),
\end{equation}
where $h_L$ and $h_R$ are the axial and radial edge-to-center density ratios.
These are computed from:
\begin{align}
  h_{L,0} &= \frac{0.86}{1+\alpha_0}\left[
    3 + \frac{\eta\,L}{2\lambda_i}
    + \left(\frac{0.86\,L\,u_B}{\pi\,D_a}\right)^{\!2}
  \right]^{-1/2}, \label{eq:hL}\\[4pt]
  h_{R,0} &= \frac{0.80}{1+\alpha_0}\left[
    4 + \frac{\eta\,R}{\lambda_i}
    + \left(\frac{0.80\,R\,u_B}{\chi_{01}\,J_1(\chi_{01})\,D_a}\right)^{\!2}
  \right]^{-1/2}, \label{eq:hR}
\end{align}
with $\chi_{01} = 2.405$, $J_1(\chi_{01}) = 0.5191$, $\eta = 2T_+/(T_+ + T_-)$,
and $\alpha_0 \approx 1.5\,\alpha$.

For electronegative plasmas ($\alpha > 0.1$), a correction factor $h_c$ arising
from ion--ion neutralization in the bulk is added in quadrature:
\begin{equation}
  \label{eq:hc}
  h_{L} = \sqrt{h_{L,0}^2 + h_c^2}, \qquad
  h_{R} = \sqrt{h_{R,0}^2 + h_c^2}.
\end{equation}

\paragraph{Ambipolar diffusion coefficient:}
\begin{equation}
  \label{eq:Da}
  D_a = D_+\;\frac{1 + \gamma_+ + \gamma_+\,\alpha_s}{1 + \gamma_+\,\alpha_s},
\end{equation}
where $D_+ = k_BT_+\lambda_i/(m_i\bar{v}_i)$ is the free ion diffusion
coefficient, $\gamma_+ = \Te/T_+$, and $\lambda_i =
1/(n_{\text{total}}\,\sigma_{in})$ with $\sigma_{in} \approx
\SI{1e-18}{m^2}$.

\paragraph{Sheath-edge electronegativity:}
The electronegativity at the sheath boundary $\alpha_s$ is obtained by solving
iteratively:
\begin{equation}
  \label{eq:alpha_s}
  \alpha = \alpha_s\;\exp\!\left[
    \frac{(1+\alpha_s)(\gamma_- - 1)}{2(1+\gamma_-\,\alpha_s)}
  \right],
\end{equation}
where $\alpha = n_-/\ne$ is the bulk electronegativity.

\paragraph{Ion neutralization at wall:}
Positive ions reaching the wall are neutralized, returning their corresponding
neutral species to the gas phase (e.g., Ar$^+ \to$ Ar, NF$_2^+ \to$ NF$_2$).


\subsection{Pumping Model}
\label{sec:pumping}

The model operates under constant-flow conditions with a fixed throughput.

\paragraph{Pump-in rate:}
Only the feed gases (Ar and NF$_3$) are injected:
\begin{equation}
  \label{eq:pump_in}
  \dot{n}_{X,\text{in}} = \frac{4.48 \times 10^{17}\;Q_X}{V},
\end{equation}
where $Q_X$ is the flow rate of species $X$ in sccm and $V$ is the chamber
volume in m$^3$. The numerical prefactor converts sccm to particles\,s$^{-1}$.

\paragraph{Pump-out rate:}
All neutral species and ions are pumped out at a rate proportional to their
density:
\begin{equation}
  \label{eq:pump_out}
  \dot{n}_{X,\text{out}} = -\frac{1.27 \times 10^{-5}\;Q_{\text{total}}}{c_{cc}\;p\;V}\;n_X,
\end{equation}
where $Q_{\text{total}}$ is the total flow rate (sccm), $p$ is the operating
pressure (Torr), $c_{cc}$ is a consistency factor relating actual total density
to the ideal gas density at the set pressure and gas temperature, and $V$ is
the volume.


\subsection{Electron Energy (Power) Balance}
\label{sec:power}

The electron temperature $\Te$ is obtained from an energy balance equation:
\begin{equation}
  \label{eq:power_balance}
  \boxed{
  \ddt{}\!\left(\frac{3}{2}\,e\,\ne\,\Te\right)
    = \frac{P_{\text{abs}}}{V}
    - \sum_r \ne\,n_r\,k_r(\Te)\,\varepsilon_r
    - \sum_i n_i\,u_{B,i}\,\frac{A_{\text{eff}}}{V}\,\varepsilon_{\text{wall}}
  }
\end{equation}
where $P_{\text{abs}}$ is the absorbed RF power, $\varepsilon_r$ is the energy
lost per collision in reaction $r$ (the threshold energy for inelastic
collisions, and $3\Te m_e/M$ for elastic collisions), and $\varepsilon_{\text{wall}}$
is the total energy lost per ion--electron pair reaching the wall:
\begin{equation}
  \label{eq:eps_wall}
  \varepsilon_{\text{wall}} = \varepsilon_e + \varepsilon_p + \varepsilon_s,
\end{equation}
with:
\begin{align}
  \varepsilon_e &= 2\Te
    &\quad&\text{(electron kinetic energy)}, \label{eq:eps_e}\\
  \varepsilon_p &= \frac{1+\alpha_s}{2+2\alpha_s\gamma_-}\;\Te
    &\quad&\text{(presheath drop)}, \label{eq:eps_p}\\
  \varepsilon_s &= \Te\ln\!\left(\frac{4u_B}{v_e}\;\frac{1+\alpha_s}{1+\alpha_s(v_-/v_e)^2}\right)
    &\quad&\text{(sheath drop)}, \label{eq:eps_s}
\end{align}
where $v_e = \sqrt{8e\Te/(\pi m_e)}$ is the electron thermal velocity.

In practice, the ODE is rewritten for $\Te$ directly:
\begin{equation}
  \dv{\Te}{t} = \frac{2}{3\,e\,\ne}\left[
    \frac{P_{\text{abs}}}{V}
    - e\sum_r \ne\,n_r\,k_r\,\varepsilon_r
    - e\sum_i n_i\,u_{B,i}\,\frac{A_{\text{eff}}}{V}\,\varepsilon_{\text{wall}}
  \right].
\end{equation}


\subsection{Gas Temperature Model}
\label{sec:Tgas}

The gas temperature is computed from an empirical correlation rather than
a separate energy balance:
\begin{align}
  \Tg &= 300 + s(p)\,\log_{10}\!\left(\frac{P_{\text{abs}}}{40}\right), \label{eq:Tg}\\
  s(p) &= 780\bigl(1 - e^{-0.091\,p}\bigr) + 250\,e^{-0.337\,p}, \label{eq:sp}
\end{align}
where $p$ is the pressure in mTorr and $P_{\text{abs}}$ in watts.

\subsection{Ion Temperature}

The ion and negative-ion temperatures are estimated from the Lee--Lieberman
empirical formula:
\begin{equation}
  \label{eq:Tion}
  T_+ = T_- = \frac{0.5\mev \times 11604.5\;\text{K} - \Tg}{p\;[\text{mTorr}]} + \Tg.
\end{equation}
At pressures above a few mTorr, $T_+ \approx \Tg$.


% ============================================================================
\section{Reactor Geometry and Operating Conditions}
\label{sec:reactor}
% ============================================================================

The model parameters correspond to the ICP reactor described by Huang
\textit{et al.}\ (2026).

\begin{table}[H]
  \centering
  \caption{Reactor geometry and base-case operating conditions.}
  \label{tab:reactor}
  \begin{tabular}{lll}
    \toprule
    \textbf{Parameter} & \textbf{Value} & \textbf{Description} \\
    \midrule
    $R$           & \SI{21}{cm}    & Chamber radius \\
    $L$           & \SI{6}{cm}     & Chamber height \\
    $V$           & \SI{8313}{cm^3} & Chamber volume ($\pi R^2 L$) \\
    $A$           & \SI{3554}{cm^2} & Total wall area \\
    $P_{\text{abs}}$ & \SIrange{500}{1000}{W} & Absorbed RF power \\
    $p$           & \SIrange{10}{100}{mTorr} & Operating pressure \\
    $Q_{\text{Ar}}$  & \SI{80}{sccm} & Argon flow rate (base case) \\
    $Q_{\text{NF}_3}$ & \SI{20}{sccm} & NF$_3$ flow rate (base case, 20\%) \\
    $T_{\text{wall}}$ & \SI{300}{K}   & Wall temperature \\
    \bottomrule
  \end{tabular}
\end{table}


% ============================================================================
\section{Numerical Implementation}
\label{sec:numerics}
% ============================================================================

\subsection{Solution Strategy}

The system of 20 coupled ODEs (19 species densities + $\Te$) is integrated in
time using \texttt{scipy.integrate.solve\_ivp} with the BDF (Backward
Differentiation Formula) method, which is well suited for stiff systems. The
integration proceeds from $t = 0$ to $t = \SI{50}{ms}$, by which time a steady
state has been reached.

\subsection{Initial Conditions}

The initial neutral densities are obtained from the ideal gas law at
$T = \SI{300}{K}$:
\begin{equation}
  n_{\text{total},0} = \frac{p}{k_B \times 300\;\text{K}},
\end{equation}
partitioned according to the feed composition (e.g., 80\% Ar, 20\% NF$_3$).
Ion densities are initialized at small seed values ($10^{14}$\,m$^{-3}$) and the
initial electron temperature is set to 4~eV.

\subsection{Constraint Handling}

The electron density is computed at each RHS evaluation from charge neutrality
(Eq.~\ref{eq:neutrality}) rather than as an independent ODE variable. All
species densities are clamped to non-negative values.

\subsection{Convergence}

Typical solutions converge in $\sim$1500 function evaluations with relative and
absolute tolerances of $10^{-8}$ and $10^{6}$, respectively.


% ============================================================================
\section{Results}
\label{sec:results}
% ============================================================================

All results presented below use the base-case conditions of
Table~\ref{tab:reactor} (80/20 Ar/NF$_3$, 30~mTorr, 600~W) unless otherwise
noted.

\subsection{Base-Case Steady State}

Table~\ref{tab:basecase} summarizes the key outputs at the base-case conditions.

\begin{table}[H]
  \centering
  \caption{Steady-state results at the base-case conditions (600~W, 30~mTorr,
  20\% NF$_3$).  The reference values are taken from Huang \textit{et al.}\
  (2026) at matching conditions.}
  \label{tab:basecase}
  \begin{tabular}{lccc}
    \toprule
    \textbf{Quantity} & \textbf{Model} & \textbf{Reference} & \textbf{Agreement} \\
    \midrule
    $\Te$                   & 2.40~eV   & $\sim$2.5~eV   & Excellent \\
    $\ne$                   & $1.2\times10^{17}$~m$^{-3}$
                            & $\sim7\times10^{16}$~m$^{-3}$ & $\sim$1.7$\times$ high \\
    $\Tg$                   & 1158~K    & ---             & Empirical formula \\
    NF$_3$ dissociation     & 98\%      & $\sim$93\%      & Good \\
    Dominant cation         & Ar$^+$ (76\%)  & Ar$^+$ (74\%)  & \checkmark \\
    Dominant anion          & F$^-$ (99.9\%) & F$^-$          & \checkmark \\
    Electronegativity $\alpha$ & 0.35   & $\sim$0.3       & Excellent \\
    \bottomrule
  \end{tabular}
\end{table}

The remaining $\sim$1.7$\times$ overestimate in $\ne$ is within the typical
accuracy range of global models, which commonly achieve factor-of-two agreement
with experiment due to uncertainties in cross sections, wall loss parameters,
and power coupling efficiency.


\subsection{Electron Density versus Absorbed Power}

Figure~\ref{fig:ne_power} shows the electron density as a function of absorbed
power at 30~mTorr with 20\% NF$_3$. The model predicts a nearly linear increase
of $\ne$ with power, rising from $\sim9.3\times10^{16}$~m$^{-3}$ at 500~W to
$\sim2.5\times10^{17}$~m$^{-3}$ at 1000~W. This linear scaling is characteristic
of ICP discharges where the electron temperature is approximately constant and
additional power is channeled into increased ionization. The trend is consistent
with the experimental cutoff-probe measurements and model results reported in the
reference paper (Fig.~10a therein).

\begin{figure}[H]
  \centering
  \includegraphics[width=0.6\textwidth]{figures/fig_ne_vs_power.pdf}
  \caption{Electron density versus absorbed power at 30~mTorr and 20\% NF$_3$
  feed composition. The near-linear dependence reflects the constant-$\Te$
  regime typical of ICP discharges. Compare with Fig.~10a of the reference.}
  \label{fig:ne_power}
\end{figure}


\subsection{Electron Density versus Pressure}

Figure~\ref{fig:ne_pressure} shows $\ne$ as a function of operating pressure at
600~W. The electron density decreases monotonically from
$\sim1.6\times10^{17}$~m$^{-3}$ at 10~mTorr to $\sim5.4\times10^{16}$~m$^{-3}$
at 100~mTorr. This decrease occurs because higher neutral density enhances
electron attachment (forming F$^-$), which consumes electrons, and because
enhanced collisionality reduces the effective ionization rate. The monotonic
decrease is qualitatively consistent with the reference model (Fig.~10b therein),
although the reference shows an initial rise at very low NF$_3$ fractions that
our single-composition sweep does not capture.

\begin{figure}[H]
  \centering
  \includegraphics[width=0.6\textwidth]{figures/fig_ne_vs_pressure.pdf}
  \caption{Electron density versus pressure at 600~W and 20\% NF$_3$ feed.
  The decrease with pressure reflects enhanced electron attachment at higher
  neutral densities. Compare with Fig.~10b of the reference.}
  \label{fig:ne_pressure}
\end{figure}


\subsection{Neutral Species Densities versus Pressure}

Figure~\ref{fig:neutrals} shows the densities of all tracked neutral species as
a function of pressure. Argon is the most abundant species at all pressures due
to its 80\% feed fraction and chemical inertness. Among the dissociation products,
fluorine atoms (F) dominate, followed by N$_2$ and F$_2$, consistent with the
highly dissociated NF$_3$ regime. The residual NF$_3$ density is orders of
magnitude below the initial feed density, confirming near-complete dissociation.
The NF$_2$ and NF fragment densities remain low ($10^{17}$--$10^{18}$~m$^{-3}$),
indicating rapid further dissociation of these intermediates. As pressure
increases, the total neutral density rises (ideal gas scaling), and the relative
abundance of F atoms increases due to the increased dissociative electron
attachment rate.

\begin{figure}[H]
  \centering
  \includegraphics[width=0.65\textwidth]{figures/fig_neutrals_vs_pressure.pdf}
  \caption{Neutral species densities versus pressure at 600~W and 20\% NF$_3$
  feed. Ar is the majority species; F is the dominant dissociation product.
  NF$_3$ is nearly fully dissociated at all pressures shown.
  Compare with Fig.~7b of the reference.}
  \label{fig:neutrals}
\end{figure}


\subsection{Electron and Gas Temperature versus Pressure}

Figure~\ref{fig:Te_Tg} presents the electron temperature $\Te$ (left axis) and
gas temperature $\Tg$ (right axis) as functions of pressure. $\Te$ decreases
from 2.58~eV at 10~mTorr to 2.15~eV at 100~mTorr, which is the standard behavior
for ICP discharges: at higher pressure, more frequent collisions lead to enhanced
energy transfer to the gas, lowering the electron temperature needed to sustain
the discharge. This range and trend are in good agreement with the reference
(Fig.~11b therein), which reports $\Te \approx 2.5$--$1.5$~eV over the same
pressure range. The gas temperature, computed from the empirical model
(Eqs.~\ref{eq:Tg}--\ref{eq:sp}), increases from $\sim$860~K at 10~mTorr to
$\sim$1220~K at 100~mTorr, reflecting enhanced power deposition at higher
pressures.

\begin{figure}[H]
  \centering
  \includegraphics[width=0.6\textwidth]{figures/fig_Te_Tg_vs_pressure.pdf}
  \caption{Electron temperature (left, blue) and gas temperature (right, red)
  versus pressure at 600~W and 20\% NF$_3$. $\Te$ decreases with pressure as
  collisional energy transfer increases, while $\Tg$ increases due to enhanced
  neutral heating. Compare with Fig.~11b of the reference.}
  \label{fig:Te_Tg}
\end{figure}


\subsection{Electronegativity versus Pressure}

Figure~\ref{fig:alpha} shows the electronegativity $\alpha = n_-/\ne$ as a
function of pressure. $\alpha$ increases from $\sim$0.15 at 10~mTorr to
$\sim$0.75 at 100~mTorr. This increase is driven by the growing importance of
dissociative electron attachment (e + NF$_3 \to$ NF$_2$ + F$^-$ and
e + NF$_2 \to$ NF + F$^-$) at higher pressures where the NF$_x$ target density
is larger. The values are in excellent agreement with the reference, which reports
$\alpha \approx 0.1$--1.2 over a similar pressure range.

\begin{figure}[H]
  \centering
  \includegraphics[width=0.6\textwidth]{figures/fig_alpha_vs_pressure.pdf}
  \caption{Electronegativity $\alpha = n_-/\ne$ versus pressure at 600~W and
  20\% NF$_3$. The increase with pressure reflects enhanced dissociative
  electron attachment. Compare with Fig.~14 of the reference.}
  \label{fig:alpha}
\end{figure}


\subsection{Ion Composition}

Figure~\ref{fig:ions} shows the positive and negative ion composition at the
base-case conditions. Ar$^+$ is the dominant positive ion (76\%), followed by
N$_2^+$ (9\%) and N$^+$ (6\%). This is in excellent agreement with the reference
model (Fig.~4b therein), which predicts Ar$^+$ dominance at $\sim$74\% for 20\%
NF$_3$ mixtures. The dominance of Ar$^+$ arises because: (i)~Ar is the majority
species, and (ii)~the ionization threshold of Ar (15.76~eV) is offset by its
high density, giving a large absolute ionization rate. Among the negative ions,
F$^-$ is overwhelmingly dominant (99.8\%), produced primarily by dissociative
electron attachment to NF$_3$ and NF$_2$.

\begin{figure}[H]
  \centering
  \includegraphics[width=0.75\textwidth]{figures/fig_ion_composition.pdf}
  \caption{Ion composition at the base-case conditions (600~W, 30~mTorr, 20\%
  NF$_3$). Left: positive ions; right: negative ions. Ar$^+$ dominates the
  cation population and F$^-$ dominates the anion population.
  Compare with Fig.~4 of the reference.}
  \label{fig:ions}
\end{figure}


\subsection{NF$_3$ Dissociation versus Pressure}

Figure~\ref{fig:dissoc} shows the NF$_3$ dissociation fraction versus pressure.
The model predicts near-complete ($>$98\%) dissociation of NF$_3$ across the
entire pressure range studied (10--100~mTorr), consistent with the reference
which reports $>$90\% dissociation at these conditions. The very high
dissociation is driven by the combination of: (i)~low-threshold dissociative
electron attachment (e + NF$_3 \to$ NF$_2$ + F$^-$, threshold $\approx$
0.1~eV), which is a thermal process highly efficient at all electron
temperatures; (ii)~electron-impact dissociation (e + NF$_3 \to$ NF$_2$ + F + e,
threshold $\approx$ 2.5~eV); and (iii)~Penning dissociation by metastable Ar
atoms (Ar$^* +$ NF$_3 \to$ products). The slight over-prediction (98\% vs.\
$\sim$93\%) compared with the reference may reflect small differences in the
effective pumping and gas residence time parameters.

\begin{figure}[H]
  \centering
  \includegraphics[width=0.6\textwidth]{figures/fig_dissociation_vs_pressure.pdf}
  \caption{NF$_3$ dissociation fraction versus pressure at 600~W and 20\%
  NF$_3$. Near-complete dissociation is predicted at all pressures, driven by
  highly efficient dissociative electron attachment.
  Compare with Fig.~5b of the reference.}
  \label{fig:dissoc}
\end{figure}


\subsection{Six-Panel Summary}

Figure~\ref{fig:summary} provides an overview of all parametric trends on a
single page for convenient reference.

\begin{figure}[H]
  \centering
  \includegraphics[width=\textwidth]{figures/fig_summary_6panel.pdf}
  \caption{Six-panel summary of model results. (a)~Electron density vs.\ power.
  (b)~Electron density vs.\ pressure. (c)~Neutral species densities vs.\
  pressure on a logarithmic scale. (d)~Electron temperature vs.\ pressure.
  (e)~Electronegativity vs.\ pressure. (f)~NF$_3$ dissociation fraction vs.\
  pressure. All panels use 20\% NF$_3$ feed; panels (a) use 30~mTorr;
  panels (b)--(f) use 600~W.}
  \label{fig:summary}
\end{figure}


% ============================================================================
\section{Discussion}
\label{sec:discussion}
% ============================================================================

\subsection{Sources of Quantitative Discrepancy}

The primary quantitative discrepancy is the $\sim$1.7$\times$ overestimate of the
electron density compared to the reference. This level of agreement is typical
for 0D global models and falls within the commonly accepted factor-of-two accuracy.
Several contributing factors remain:

\begin{enumerate}[nosep]
  \item \textbf{Ion--neutral collision cross section.} The ion mean free path
    $\lambda_i$ depends on $\sigma_{in}$, which was set to
    $\SI{5e-19}{m^2}$. This value directly affects the $h$-factors
    (Eqs.~\ref{eq:hL}--\ref{eq:hR}) and thus the ion wall-loss rate.
  \item \textbf{Electronegative correction ($h_c$).} The Thorsteinsson--Gudmundsson
    $h_c$ correction involves several parameters (ion--ion recombination rate,
    critical density $n^*$) that are not precisely known and were estimated.
  \item \textbf{Power transfer efficiency.} The model assumes 100\% coupling of
    $P_{\text{abs}}$ to electrons. In practice, some power may be absorbed by
    ions or dissipated in the matching network.
  \item \textbf{Gas temperature model.} The empirical $\Tg$ formula
    (Eq.~\ref{eq:Tg}) was not specifically calibrated for this reactor geometry.
\end{enumerate}


\subsection{Critical Implementation Issues Identified}

Two critical implementation errors were identified and corrected during model
development:

\paragraph{Unit conversion in N$_2$ rate coefficients (R54--R57).}
These rate coefficients are parameterized with the electron temperature in
\textit{Kelvin}, not eV. Since the power-law prefactor $\Te^B$ is exquisitely
sensitive to the choice of units (a factor of $\sim$40{,}000 for $B = -1.13$),
incorrect conversion leads to a catastrophic overestimate of the N$_2$(C)
excitation rate, which then dominates the power balance and starves all other
channels of energy. Before this fix, the electron density was $\sim$70$\times$
too low.

\paragraph{Misidentified reactant in R86.}
Reaction R86 was initially entered as e~+~F$_2 \to$~F~+~F~+~2e (dissociative
ionization of neutral F$_2$), when the correct reaction is
e~+~F$_2^- \to$~F~+~F~+~e (dissociative detachment of the negative ion
F$_2^-$). The erroneous reaction created a spurious ionization channel from the
abundant neutral F$_2$ population, which accounted for 73\% of all electron
production in the model. After correction, $\ne$ decreased by a factor of
$\sim$2.3 and the electronegativity $\alpha$ increased from 0.06 to 0.35,
bringing it into excellent agreement with the reference value of $\sim$0.3.

These examples highlight the sensitivity of global plasma models to individual
rate coefficients and underscore the importance of careful reaction-by-reaction
verification against the source literature.


\subsection{Dominant Reaction Pathways}

The model reveals the following dominant pathways at the base-case conditions:

\paragraph{NF$_3$ dissociation:} Dominated by dissociative electron attachment
(R4: $\ee + \text{NF}_3 \to \text{NF}_2 + \text{F}^-$) and electron-impact
dissociation (R3: $\ee + \text{NF}_3 \to \text{NF}_2 + \text{F} + \ee$). The
attachment channel is particularly efficient due to its near-zero threshold
energy.

\paragraph{Electron production:} Primarily from Ar ionization (R67: $\ee +
\text{Ar} \to \text{Ar}^+ + 2\ee$), with secondary contributions from NF$_2$
and NF ionization.

\paragraph{Electron loss:} Dissociative attachment to NF$_3$, NF$_2$, and F$_2$,
plus wall losses via ambipolar diffusion.

\paragraph{F atom production:} From electron-impact dissociation of NF$_3$,
NF$_2$, NF, and F$_2$, as well as from dissociative attachment.

\paragraph{N$_2$ formation:} Primarily through neutral--neutral recombination
(R94: N + NF $\to$ N$_2$ + F) and wall recombination (N + N(wall) $\to$
N$_2$).


% ============================================================================
\section{Summary and Conclusions}
\label{sec:conclusions}
% ============================================================================

A complete 0D global model for Ar/NF$_3$ ICP discharges has been formulated,
implemented in Python, and validated against the published results of Huang
\textit{et al.}\ (2026). The model includes:

\begin{itemize}[nosep]
  \item 21 plasma species (10 neutrals, 8 positive ions, 2 negative ions, electrons).
  \item 155 volumetric reactions spanning electron-impact, Penning, neutral,
    ion--neutral, ion--ion, thermal, and three-body processes.
  \item Surface recombination with species-dependent sticking coefficients.
  \item A self-consistent electron energy balance with collisional and wall
    energy losses.
  \item The Thorsteinsson--Gudmundsson ion wall-loss model for electronegative
    plasmas.
  \item An empirical gas temperature model.
\end{itemize}

The model correctly reproduces all key qualitative trends:
\begin{itemize}[nosep]
  \item Linear increase of $\ne$ with absorbed power.
  \item Decrease of $\ne$ and $\Te$ with increasing pressure.
  \item Near-complete dissociation of NF$_3$ ($>$98\%).
  \item Ar$^+$ as the dominant positive ion (76\%, reference: 74\%) and F$^-$ as
    the dominant negative ion.
  \item Increasing electronegativity with pressure ($\alpha \approx 0.35$ at
    30~mTorr, reference: $\sim$0.3).
  \item Electron temperature $\Te \approx 2.4$~eV at the base case (reference:
    $\sim$2.5~eV).
\end{itemize}

The absolute electron density agrees with the reference within a factor of 1.7,
which is within the typical accuracy of zero-dimensional global models.
The remaining discrepancy is attributed to uncertainties in the ion wall-loss
parameters and power coupling efficiency.


% ============================================================================
% Appendix: Physical Constants and Nomenclature
% ============================================================================
\appendix
\section{Physical Constants}

\begin{table}[H]
  \centering
  \begin{tabular}{lll}
    \toprule
    \textbf{Symbol} & \textbf{Value} & \textbf{Description} \\
    \midrule
    $e$           & $1.602\times10^{-19}$~C     & Elementary charge \\
    $k_B$         & $1.381\times10^{-23}$~J/K   & Boltzmann constant \\
    $m_e$         & $9.109\times10^{-31}$~kg    & Electron mass \\
    $\varepsilon_0$ & $8.854\times10^{-12}$~F/m & Vacuum permittivity \\
    1~eV          & $11604.5$~K                  & Temperature equivalent \\
    \bottomrule
  \end{tabular}
\end{table}

\section{Surface Recombination Probabilities}

\begin{table}[H]
  \centering
  \begin{tabular}{lc}
    \toprule
    \textbf{Species} & $\gamma$ \\
    \midrule
    F        & 0.01 \\
    N        & 0.07 \\
    NF       & 0.01 \\
    NF$_2$   & 0.01 \\
    NF$_3$   & 0.0  \\
    Ar(1$s_5$) & 1.0 \\
    Ar(4$p$)   & 1.0 \\
    \bottomrule
  \end{tabular}
\end{table}


\end{document}
